\subsection{The cycle decomposition}

Next, we shall discuss the \emph{cycle decomposition} (or \emph{disjoint cycle
decomposition}) of a permutation. Again, this is a fairly well-known and
elementary tool, so we will restrict ourselves to the basic properties and
omit the details.

We begin with an introductory example:

\begin{example}
\label{exa.perm.dcd.1}Let $\sigma\in S_{9}$ be the permutation with one-line
notation $461352987$. Here is its cycle digraph:%
\[%
% [tikz figure removed]
%BeginExpansion
% [tikz figure removed]%
%EndExpansion
\qquad\qquad%
% [tikz figure removed]
%BeginExpansion
% [tikz figure removed]%
%EndExpansion
\qquad%
%TCIMACRO{\TeXButton{tikz cycle 5}{\raisebox{8ex}{
%% [tikz figure removed]
%}}}%
%BeginExpansion
\raisebox{8ex}{
% [tikz figure removed]
}%
%EndExpansion
\qquad%
% [tikz figure removed]
%BeginExpansion
% [tikz figure removed]%
%EndExpansion
\qquad%
%TCIMACRO{\TeXButton{tikz cycle 8}{\raisebox{8ex}{
%% [tikz figure removed]
%}}}%
%BeginExpansion
\raisebox{8ex}{
% [tikz figure removed]
}%
%EndExpansion
\]
(where we have strategically arranged the cycles apart horizontally). This
digraph consists of five node-disjoint cycles (i.e., cycles that share no
nodes); thus, we can view the permutation $\sigma$ as acting on the five
subsets $\left\{  1,4,3\right\}  $, $\left\{  2,6\right\}  $, $\left\{
5\right\}  $, $\left\{  7,9\right\}  $ and $\left\{  8\right\}  $ of $\left[
9\right]  $ separately. On the first of these five subsets, $\sigma$ acts as
the $3$-cycle $\operatorname*{cyc}\nolimits_{1,4,3}$ (in the sense that
$\sigma\left(  k\right)  =\operatorname*{cyc}\nolimits_{1,4,3}\left(
k\right)  $ for each $k\in\left\{  1,4,3\right\}  $). On the second, it acts
as the $2$-cycle $\operatorname*{cyc}\nolimits_{2,6}$. On the third, it acts
as the $1$-cycle $\operatorname*{cyc}\nolimits_{5}$ (which, of course, is the
identity map). On the fourth, it acts as the $2$-cycle $\operatorname*{cyc}%
\nolimits_{7,9}$. On the fifth, it acts as the $1$-cycle $\operatorname*{cyc}%
\nolimits_{8}$ (which, again, is just the identity map). Combining these
observations, we conclude that%
\[
\sigma\left(  k\right)  =\left(  \operatorname*{cyc}\nolimits_{1,4,3}%
\circ\operatorname*{cyc}\nolimits_{2,6}\circ\operatorname*{cyc}\nolimits_{5}%
\circ\operatorname*{cyc}\nolimits_{7,9}\circ\operatorname*{cyc}\nolimits_{8}%
\right)  \left(  k\right)  \ \ \ \ \ \ \ \ \ \ \text{for each }k\in\left[
9\right]
\]
(because, when we apply the composed permutation $\operatorname*{cyc}%
\nolimits_{1,4,3}\circ\operatorname*{cyc}\nolimits_{2,6}\circ
\operatorname*{cyc}\nolimits_{5}\circ\operatorname*{cyc}\nolimits_{7,9}%
\circ\operatorname*{cyc}\nolimits_{8}$ to an element of $\left[  9\right]  $,
then four of the five cycles $\operatorname*{cyc}\nolimits_{1,4,3}%
,\operatorname*{cyc}\nolimits_{2,6},\operatorname*{cyc}\nolimits_{5}%
,\operatorname*{cyc}\nolimits_{7,9},\operatorname*{cyc}\nolimits_{8}$ will
leave this element unchanged, whereas the remaining one will move it one step
forward along the appropriate cycle of the above digraph -- which is precisely
what the permutation $\sigma$ does to our element). This entails that%
\begin{equation}
\sigma=\operatorname*{cyc}\nolimits_{1,4,3}\circ\operatorname*{cyc}%
\nolimits_{2,6}\circ\operatorname*{cyc}\nolimits_{5}\circ\operatorname*{cyc}%
\nolimits_{7,9}\circ\operatorname*{cyc}\nolimits_{8}.
\label{eq.exa.perm.dcd.1.1}%
\end{equation}

\end{example}

In Example \ref{exa.perm.dcd.1}, we have represented our permutation
$\sigma\in S_{9}$ as a composition of five cycles with the property that each
element of $\left[  9\right]  $ appears in exactly one of these cycles. This
is not specific to the permutation $\sigma$ chosen in Example
\ref{exa.perm.dcd.1}. Indeed, for any finite set $X$, any permutation
$\sigma\in S_{X}$ can be written as a composition of finitely many cycles
$\operatorname*{cyc}\nolimits_{i_{1},i_{2},\ldots,i_{k}}$ with the property
that each element of $X$ appears in exactly one of these cycles. Moreover,
this representation of $\sigma$ is unique up to

\begin{itemize}
\item swapping the cycles (for example, we could have replaced
$\operatorname*{cyc}\nolimits_{1,4,3}\circ\operatorname*{cyc}\nolimits_{2,6}$
by $\operatorname*{cyc}\nolimits_{2,6}\circ\operatorname*{cyc}%
\nolimits_{1,4,3}$ in (\ref{eq.exa.perm.dcd.1.1})), and

\item rotating each cycle (for example, we could have replaced
$\operatorname*{cyc}\nolimits_{1,4,3}$ by $\operatorname*{cyc}%
\nolimits_{4,3,1}$ or by $\operatorname*{cyc}\nolimits_{3,1,4}$ in
(\ref{eq.exa.perm.dcd.1.1})).
\end{itemize}

Let us state this in a more rigorous fashion:

\begin{theorem}
[disjoint cycle decomposition of permutations]\label{thm.perm.dcd.main}Let $X$
be a finite set. Let $\sigma$ be a permutation of $X$. Then: \medskip

\textbf{(a)} There is a list%
\begin{align*}
&  \Big(\left(  a_{1,1},a_{1,2},\ldots,a_{1,n_{1}}\right)  ,\\
&  \ \ \ \left(  a_{2,1},a_{2,2},\ldots,a_{2,n_{2}}\right)  ,\\
&  \ \ \ \ldots,\\
&  \ \ \ \left(  a_{k,1},a_{k,2},\ldots,a_{k,n_{k}}\right)  \Big)
\end{align*}
of nonempty lists of elements of $X$ such that:

\begin{itemize}
\item each element of $X$ appears exactly once in the composite list%
\begin{align*}
&  (a_{1,1},a_{1,2},\ldots,a_{1,n_{1}},\\
&  \ \ \ a_{2,1},a_{2,2},\ldots,a_{2,n_{2}},\\
&  \ \ \ \ldots,\\
&  \ \ \ a_{k,1},a_{k,2},\ldots,a_{k,n_{k}}),
\end{align*}
and

\item we have%
\[
\sigma=\operatorname*{cyc}\nolimits_{a_{1,1},a_{1,2},\ldots,a_{1,n_{1}}}%
\circ\operatorname*{cyc}\nolimits_{a_{2,1},a_{2,2},\ldots,a_{2,n_{2}}}%
\circ\cdots\circ\operatorname*{cyc}\nolimits_{a_{k,1},a_{k,2},\ldots
,a_{k,n_{k}}}.
\]

\end{itemize}

Such a list is called a \emph{disjoint cycle decomposition} (or short
\emph{DCD}) of $\sigma$. Its entries (which themselves are lists of elements
of $X$) are called the \emph{cycles} of $\sigma$. \medskip

\textbf{(b)} Any two DCDs of $\sigma$ can be obtained from each other by
(repeatedly) swapping the cycles with each other, and rotating each cycle
(i.e., replacing $\left(  a_{i,1},a_{i,2},\ldots,a_{i,n_{i}}\right)  $ by
$\left(  a_{i,2},a_{i,3},\ldots,a_{i,n_{i}},a_{i,1}\right)  $). \medskip

\textbf{(c)} Now assume that $X$ is a set of integers (or, more generally, any
totally ordered finite set). Then, there is a unique DCD
\begin{align*}
&  \Big(\left(  a_{1,1},a_{1,2},\ldots,a_{1,n_{1}}\right)  ,\\
&  \ \ \ \left(  a_{2,1},a_{2,2},\ldots,a_{2,n_{2}}\right)  ,\\
&  \ \ \ \ldots,\\
&  \ \ \ \left(  a_{k,1},a_{k,2},\ldots,a_{k,n_{k}}\right)  \Big)
\end{align*}
of $\sigma$ that satisfies the additional requirements that

\begin{itemize}
\item we have $a_{i,1}\leq a_{i,p}$ for each $i\in\left[  k\right]  $ and each
$p\in\left[  n_{i}\right]  $ (that is, each cycle in this DCD is written with
its smallest entry first), and

\item we have $a_{1,1}>a_{2,1}>\cdots>a_{k,1}$ (that is, the cycles appear in
this DCD in the order of decreasing first entries).
\end{itemize}
\end{theorem}

\begin{example}
Let $\sigma\in S_{9}$ be the permutation from Example \ref{exa.perm.dcd.1}.
Then, the representation (\ref{eq.exa.perm.dcd.1.1}) shows that%
\[
\Big(\left(  1,4,3\right)  ,\ \left(  2,6\right)  ,\ \left(  5\right)
,\ \left(  7,9\right)  ,\ \left(  8\right)  \Big)
\]
is a DCD of $\sigma$. By swapping the five cycles of this DCD, and by rotating
each cycle, we can produce various other DCDs of $\sigma$, such as
\[
\Big(\left(  7,9\right)  ,\ \left(  6,2\right)  ,\ \left(  3,1,4\right)
,\ \left(  8\right)  ,\ \left(  5\right)  \Big).
\]
The unique DCD of $\sigma$ that satisfies the two additional requirements of
Theorem \ref{thm.perm.dcd.main} \textbf{(c)} is
\[
\Big(\left(  8\right)  ,\ \left(  7,9\right)  ,\ \left(  5\right)  ,\ \left(
2,6\right)  ,\ \left(  1,4,3\right)  \Big).
\]

\end{example}

\begin{proof}
[Proof of Theorem \ref{thm.perm.dcd.main} (sketched).]This is a classical
result with an easy proof; sadly, this easy proof does not present well in
writing. I will try to be as clear as the situation allows. Some familiarity
with digraphs (= directed graphs) is recommended\footnote{See, e.g.,
\cite[\S 5.11]{Guicha20} or \cite[\S 3.5]{Loehr-BC} for brief introductions to
digraphs.}. \medskip

\textbf{(a)} Let $\mathcal{D}$ be the cycle digraph of $\sigma$, as in Example
\ref{exa.perm.dcd.1}. This cycle digraph $\mathcal{D}$ has the following two properties:

\begin{itemize}
\item \textit{Outbound uniqueness:} For each node $i$, there is exactly one
arc outgoing from $i$. (Indeed, this is the arc from $i$ to $\sigma\left(
i\right)  $, as should be clear from the construction of $\mathcal{D}$.)

\item \textit{Inbound uniqueness:} For each node $i$, there is exactly one arc
incoming into $i$. (Indeed, this is the arc from $\sigma^{-1}\left(  i\right)
$ to $i$, since $\sigma$ is a permutation and therefore invertible.)
\end{itemize}

Using these two properties, we will now show that the cycle digraph
$\mathcal{D}$ consists of several node-disjoint cycles (i.e., several cycles
that pairwise share no nodes).

Indeed, let us first observe the following: If two cycles $C$ and $D$ of
$\mathcal{D}$ have a node in common, then they are identical\footnote{Here, we
identify any cycle with its cyclic rotations. For example, if $a\rightarrow
b\rightarrow c\rightarrow a$ is a cycle, then we consider $b\rightarrow
c\rightarrow a\rightarrow b$ to be the same cycle.} (because the outbound
uniqueness property prevents these cycles from ever separating after meeting
at the common node). In other words, any two cycles $C$ and $D$ of
$\mathcal{D}$ are either identical or node-disjoint (i.e., share no nodes with
each other).

Now, let $i$ be any node of $\mathcal{D}$. Then, if we start at $i$ and follow
the outgoing arcs, then we obtain an infinite walk%
\[
\sigma^{0}\left(  i\right)  \rightarrow\sigma^{1}\left(  i\right)
\rightarrow\sigma^{2}\left(  i\right)  \rightarrow\sigma^{3}\left(  i\right)
\rightarrow\cdots
\]
along our digraph $\mathcal{D}$. Since $X$ is finite, the Pigeonhole Principle
guarantees that this walk will eventually revisit a node it has already been
to; i.e., there exist two integers $u,v\in\mathbb{N}$ with $u<v$ and
$\sigma^{u}\left(  i\right)  =\sigma^{v}\left(  i\right)  $. Let us pick two
such integers $u,v$ with the \textbf{smallest} possible $v$. Thus,
\begin{equation}
\text{the }v\text{ nodes }\sigma^{0}\left(  i\right)  ,\sigma^{1}\left(
i\right)  ,\ldots,\sigma^{v-1}\left(  i\right)  \text{ are distinct}
\label{pf.thm.perm.dcd.main.a.dist}%
\end{equation}
(since otherwise, $v$ would not be smallest possible). Now, $\sigma$ is a
permutation and thus has an inverse $\sigma^{-1}$. Applying the map
$\sigma^{-1}$ to both sides of the equality $\sigma^{u}\left(  i\right)
=\sigma^{v}\left(  i\right)  $, we obtain $\sigma^{u-1}\left(  i\right)
=\sigma^{v-1}\left(  i\right)  $. However, if we had $u\geq1$, then
(\ref{pf.thm.perm.dcd.main.a.dist}) would entail $\sigma^{u-1}\left(
i\right)  \neq\sigma^{v-1}\left(  i\right)  $ (because $0\leq\underbrace{u}%
_{<v}-\,1<v-1$), which would contradict $\sigma^{u-1}\left(  i\right)
=\sigma^{v-1}\left(  i\right)  $. Thus, we cannot have $u\geq1$. Hence, $u<1$,
so that $u=0$. Therefore, $\sigma^{u}\left(  i\right)  =\sigma^{0}\left(
i\right)  $, so that $\sigma^{0}\left(  i\right)  =\sigma^{u}\left(  i\right)
=\sigma^{v}\left(  i\right)  $. This shows that our walk $\sigma^{0}\left(
i\right)  \rightarrow\sigma^{1}\left(  i\right)  \rightarrow\sigma^{2}\left(
i\right)  \rightarrow\sigma^{3}\left(  i\right)  \rightarrow\cdots$ is
circular: it comes back to its starting node $\sigma^{0}\left(  i\right)  =i$
after $v$ steps. We have thus found a cycle in our digraph $\mathcal{D}$:%
\[
\sigma^{0}\left(  i\right)  \rightarrow\sigma^{1}\left(  i\right)
\rightarrow\sigma^{2}\left(  i\right)  \rightarrow\cdots\rightarrow\sigma
^{v}\left(  i\right)  =\sigma^{0}\left(  i\right)  .
\]
(This is indeed a cycle, since (\ref{pf.thm.perm.dcd.main.a.dist}) shows that
its first $v$ nodes are distinct.) This shows that the node $i$ lies on a
cycle $C_{i}$ of $\mathcal{D}$ (namely, the cycle that we just found).

Now, forget that we fixed $i$. We thus have shown that each node $i$ of
$\mathcal{D}$ lies on a cycle $C_{i}$. The cycles $C_{i}$ for all nodes $i\in
X$ will be called the \emph{chosen cycles}.

Any arc of our digraph $\mathcal{D}$ must belong to one of these chosen
cycles. Indeed, if $a$ is an arc from a node $i$ to a node $j$, then $a$ must
be the only arc outgoing from $i$ (by the outbound uniqueness property); but
this means that this arc $a$ belongs to the chosen cycle $C_{i}$.

Now, let us look back at what we have shown:

\begin{itemize}
\item Any node $i$ of $\mathcal{D}$ lies on one of our chosen cycles (namely,
on $C_{i}$).

\item Some of the chosen cycles may be identical, but apart from that, the
chosen cycles are pairwise node-disjoint (since any two cycles of
$\mathcal{D}$ are either identical or node-disjoint).

\item Any arc of $\mathcal{D}$ must belong to one of these chosen cycles.
\end{itemize}

Combining these facts, we conclude that $\mathcal{D}$ consists of several
node-disjoint cycles. Let us label these cycles as%
\begin{align*}
&  \left(  a_{1,1}\rightarrow a_{1,2}\rightarrow\cdots\rightarrow a_{1,n_{1}%
}\rightarrow a_{1,1}\right)  ,\\
&  \left(  a_{2,1}\rightarrow a_{2,2}\rightarrow\cdots\rightarrow a_{2,n_{2}%
}\rightarrow a_{2,1}\right)  ,\\
&  \ldots,\\
&  \left(  a_{k,1}\rightarrow a_{k,2}\rightarrow\cdots\rightarrow a_{k,n_{k}%
}\rightarrow a_{k,1}\right)
\end{align*}
(making sure to label each cycle only once). Then, each element of $X$ appears
exactly once in the composite list%
\begin{align*}
&  (a_{1,1},a_{1,2},\ldots,a_{1,n_{1}},\\
&  \ \ \ a_{2,1},a_{2,2},\ldots,a_{2,n_{2}},\\
&  \ \ \ \ldots,\\
&  \ \ \ a_{k,1},a_{k,2},\ldots,a_{k,n_{k}}),
\end{align*}
and we have%
\[
\sigma=\operatorname*{cyc}\nolimits_{a_{1,1},a_{1,2},\ldots,a_{1,n_{1}}}%
\circ\operatorname*{cyc}\nolimits_{a_{2,1},a_{2,2},\ldots,a_{2,n_{2}}}%
\circ\cdots\circ\operatorname*{cyc}\nolimits_{a_{k,1},a_{k,2},\ldots
,a_{k,n_{k}}}%
\]
(since $\sigma$ moves any node $i\in X$ one step forward along its chosen
cycle). This proves Theorem \ref{thm.perm.dcd.main} \textbf{(a)}.

Alternative proofs of Theorem \ref{thm.perm.dcd.main} \textbf{(a)} can be
found (e.g.) in \cite[Theorem 1.5.3]{Goodman} or in \cite[\S I.4, Proposition
1.21]{Knapp16} or in \cite[Chapter I, \S 5.7, Proposition 7]{Bourba74} or in
\cite[\S 1.9, proof of Theorem 1.5.1]{Sagan19} (this is essentially our proof)
or in
\url{https://proofwiki.org/wiki/Existence_and_Uniqueness_of_Cycle_Decomposition}
(see also \cite[Exercise 7 \textbf{(e)} and \textbf{(d)}]{17f-hw7s} for a
rather formalized proof). Note that some of these sources work with a slightly
modified concept of a DCD, in which they throw away the $1$-cycles (i.e., they
replace \textquotedblleft appears exactly once\textquotedblright\ by
\textquotedblleft appears at most once\textquotedblright, and require all
cycle lengths $n_{1},n_{2},\ldots,n_{k}$ to be $>1$). For instance, the DCD
(\ref{eq.exa.perm.dcd.1.1}) becomes%
\[
\sigma=\operatorname*{cyc}\nolimits_{1,4,3}\circ\operatorname*{cyc}%
\nolimits_{2,6}\circ\operatorname*{cyc}\nolimits_{7,9}%
\]
if we use this modified notion of a DCD. \medskip

\textbf{(b)} See \cite[Theorem 1.5.3]{Goodman} or \cite[Chapter I, \S 5.7,
Proposition 7]{Bourba74}. The idea is fairly simple: Let%
\begin{align*}
&  \Big(\left(  a_{1,1},a_{1,2},\ldots,a_{1,n_{1}}\right)  ,\\
&  \ \ \ \left(  a_{2,1},a_{2,2},\ldots,a_{2,n_{2}}\right)  ,\\
&  \ \ \ \ldots,\\
&  \ \ \ \left(  a_{k,1},a_{k,2},\ldots,a_{k,n_{k}}\right)  \Big)
\end{align*}
be a DCD of $\sigma$. Then, for each $i\in X$, the cycle of this DCD that
contains $i$ is uniquely determined by $\sigma$ and $i$ up to cyclic rotation
(indeed, it is a rotated version of the list $\left(  i,\sigma\left(
i\right)  ,\sigma^{2}\left(  i\right)  ,\ldots,\sigma^{r-1}\left(  i\right)
\right)  $, where $r$ is the smallest positive integer satisfying $\sigma
^{r}\left(  i\right)  =i$). Therefore, all cycles of this DCD are uniquely
determined by $\sigma$ up to cyclic rotation and up to the relative order in
which these cycles appear in the DCD. But this is precisely the claim of
Theorem \ref{thm.perm.dcd.main} \textbf{(b)}. \medskip

\textbf{(c)} In order to obtain a DCD of $\sigma$ that satisfies these two
requirements, it suffices to

\begin{itemize}
\item start with an arbitrary DCD of $\sigma$,

\item then rotate each cycle of this DCD so that it begins with its smallest
entry, and

\item then repeatedly swap these cycles so they appear in the order of
decreasing first entries.
\end{itemize}

It is clear that the result of this procedure is uniquely determined (a
consequence of Theorem \ref{thm.perm.dcd.main} \textbf{(b)}). Thus, Theorem
\ref{thm.perm.dcd.main} \textbf{(c)} is proven.
\end{proof}

\begin{definition}
\label{def.perm.cycs.cycs}Let $X$ be a finite set. Let $\sigma$ be a
permutation of $X$. \medskip

\textbf{(a)} The \emph{cycles} of $\sigma$ are defined to be the cycles in the
DCD of $\sigma$ (as defined in Theorem \ref{thm.perm.dcd.main} \textbf{(a)}).
(This includes $1$-cycles, if there are any in the DCD of $\sigma$.)

We shall equate a cycle of $\sigma$ with any of its cyclic rotations; thus,
for example, $\left(  3,1,4\right)  $ and $\left(  1,4,3\right)  $ shall be
regarded as being the same cycle (but $\left(  3,1,4\right)  $ and $\left(
3,4,1\right)  $ shall not). \medskip

\textbf{(b)} The \emph{cycle lengths partition} of $\sigma$ shall denote the
partition of $\left\vert X\right\vert $ obtained by writing down the lengths
of the cycles of $\sigma$ in weakly decreasing order.
\end{definition}

\begin{example}
Let $\sigma\in S_{9}$ be the permutation from Example \ref{exa.perm.dcd.1}.
Then, the cycles of $\sigma$ are%
\[
\left(  1,4,3\right)  ,\ \left(  2,6\right)  ,\ \left(  5\right)  ,\ \left(
7,9\right)  ,\ \left(  8\right)  .
\]
Their lengths are $3,2,1,2,1$. Hence, the cycle lengths partition of $\sigma$
is $\left(  3,2,2,1,1\right)  $.
\end{example}

The following is obvious:

\begin{proposition}
\label{prop.perm.cycs.same}Let $X$ be a finite set. Let $i,j\in X$. Let
$\sigma$ be a permutation of $X$. Then, we have the following chain of
equivalences:%
\begin{align*}
\  &  \left(  i\text{ and }j\text{ belong to the same cycle of }\sigma\right)
\\
\Longleftrightarrow\  &  \left(  i=\sigma^{p}\left(  j\right)  \text{ for some
}p\in\mathbb{N}\right) \\
\Longleftrightarrow\  &  \left(  j=\sigma^{p}\left(  i\right)  \text{ for some
}p\in\mathbb{N}\right)  .
\end{align*}

\end{proposition}

The number of cycles of a permutation determines its sign. Let us state this
for permutations of $\left[  n\right]  $ in particular (the reader can easily
extend this to the general case using Proposition \ref{prop.perm.sign.X}):

\begin{proposition}
\label{prop.perm.cycs.sign}Let $n\in\mathbb{N}$. Let $\sigma\in S_{n}$. Let
$k\in\mathbb{N}$ be such that $\sigma$ has exactly $k$ cycles (including the
$1$-cycles). Then, $\left(  -1\right)  ^{\sigma}=\left(  -1\right)  ^{n-k}$.
\end{proposition}

\begin{proof}
[Proof of Proposition \ref{prop.perm.cycs.sign} (sketched).]Let
\begin{align*}
&  \left(  a_{1,1},a_{1,2},\ldots,a_{1,n_{1}}\right)  ,\\
&  \left(  a_{2,1},a_{2,2},\ldots,a_{2,n_{2}}\right)  ,\\
&  \ldots,\\
&  \left(  a_{k,1},a_{k,2},\ldots,a_{k,n_{k}}\right)
\end{align*}
be the $k$ cycles of $\sigma$. Thus,
\begin{align}
&  \Big(\left(  a_{1,1},a_{1,2},\ldots,a_{1,n_{1}}\right)  ,\nonumber\\
&  \ \ \ \left(  a_{2,1},a_{2,2},\ldots,a_{2,n_{2}}\right)  ,\nonumber\\
&  \ \ \ \ldots,\nonumber\\
&  \ \ \ \left(  a_{k,1},a_{k,2},\ldots,a_{k,n_{k}}\right)
\Big) \label{pf.prop.perm.cycs.sign.3}%
\end{align}
is the DCD of $\sigma$. Therefore,
\begin{align*}
\sigma &  =\operatorname*{cyc}\nolimits_{a_{1,1},a_{1,2},\ldots,a_{1,n_{1}}%
}\circ\operatorname*{cyc}\nolimits_{a_{2,1},a_{2,2},\ldots,a_{2,n_{2}}}%
\circ\cdots\circ\operatorname*{cyc}\nolimits_{a_{k,1},a_{k,2},\ldots
,a_{k,n_{k}}}\\
&  =\left(  \text{an }n_{1}\text{-cycle}\right)  \circ\left(  \text{an }%
n_{2}\text{-cycle}\right)  \circ\cdots\circ\left(  \text{an }n_{k}%
\text{-cycle}\right)  ,
\end{align*}
so that%
\begin{align}
\left(  -1\right)  ^{\sigma}  &  =\left(  -1\right)  ^{\left(  \text{an }%
n_{1}\text{-cycle}\right)  \circ\left(  \text{an }n_{2}\text{-cycle}\right)
\circ\cdots\circ\left(  \text{an }n_{k}\text{-cycle}\right)  }\nonumber\\
&  =\left(  -1\right)  ^{\left(  \text{an }n_{1}\text{-cycle}\right)  }%
\cdot\left(  -1\right)  ^{\left(  \text{an }n_{2}\text{-cycle}\right)  }%
\cdot\cdots\cdot\left(  -1\right)  ^{\left(  \text{an }n_{k}\text{-cycle}%
\right)  }\nonumber\\
&  \ \ \ \ \ \ \ \ \ \ \ \ \ \ \ \ \ \ \ \ \left(  \text{by Proposition
\ref{prop.perm.sign.props} \textbf{(e)}}\right) \nonumber\\
&  =\left(  -1\right)  ^{n_{1}-1}\cdot\left(  -1\right)  ^{n_{2}-1}\cdot
\cdots\cdot\left(  -1\right)  ^{n_{k}-1}\nonumber\\
&  \ \ \ \ \ \ \ \ \ \ \ \ \ \ \ \ \ \ \ \ \left(
\begin{array}
[c]{c}%
\text{since Proposition \ref{prop.perm.sign.props} \textbf{(c)} yields
}\left(  -1\right)  ^{\left(  \text{a }p\text{-cycle}\right)  }=\left(
-1\right)  ^{p-1}\\
\text{for any }p>0
\end{array}
\right) \nonumber\\
&  =\left(  -1\right)  ^{\left(  n_{1}-1\right)  +\left(  n_{2}-1\right)
+\cdots+\left(  n_{k}-1\right)  }. \label{pf.prop.perm.cycs.sign.5}%
\end{align}
However, recall that (\ref{pf.prop.perm.cycs.sign.3}) is a DCD of $\sigma$.
Thus, each element of $\left[  n\right]  $ appears exactly once in the
composite list%
\begin{align*}
&  (a_{1,1},a_{1,2},\ldots,a_{1,n_{1}},\\
&  \ \ \ a_{2,1},a_{2,2},\ldots,a_{2,n_{2}},\\
&  \ \ \ \ldots,\\
&  \ \ \ a_{k,1},a_{k,2},\ldots,a_{k,n_{k}}).
\end{align*}
Therefore, the length $n_{1}+n_{2}+\cdots+n_{k}$ of this composite list equals
the size $\left\vert \left[  n\right]  \right\vert =n$ of the set $\left[
n\right]  $. In other words, $n_{1}+n_{2}+\cdots+n_{k}=n$. Hence,%
\[
\left(  n_{1}-1\right)  +\left(  n_{2}-1\right)  +\cdots+\left(
n_{k}-1\right)  =\underbrace{\left(  n_{1}+n_{2}+\cdots+n_{k}\right)  }%
_{=n}-\,k=n-k.
\]
Thus, (\ref{pf.prop.perm.cycs.sign.5}) rewrites as $\left(  -1\right)
^{\sigma}=\left(  -1\right)  ^{n-k}$. This proves Proposition
\ref{prop.perm.cycs.sign}.
\end{proof}
